../cictikz/src/cictikz/data/tex/cictikz_preamble.tex

% QPSK: four symbols, one per quadrant, at +/-1 +/-j.
%
% Each is labelled twice, in cartesian and in polar form, because the
% chapter's point is that the same symbol is easier to think about in
% whichever form suits the question: cartesian when you are building the
% modulator out of an I and a Q branch, polar when you are asking what
% the phase does. The letters A to D are the symbol names the text uses.

\begin{tikzpicture}[thick]
  ../cictikz/src/cictikz/data/tex/constellation_lib.tex

  \constframe{1.6}
  \constcircle{1.414}

  \constsymbol{1}{1}
  \constsymbol{-1}{1}
  \constsymbol{-1}{-1}
  \constsymbol{1}{-1}

  \node[anchor=south west, red, scale=0.8, align=left] at (1.15,1.05)
    {$B = 1+j$\\ $= \sqrt{2}e^{j\pi/4}$};
  \node[anchor=south east, red, scale=0.8, align=right] at (-1.15,1.05)
    {$A = -1+j$\\ $= \sqrt{2}e^{j3\pi/4}$};
  \node[anchor=north east, red, scale=0.8, align=right] at (-1.15,-1.05)
    {$D = -1-j$\\ $= \sqrt{2}e^{-j3\pi/4}$};
  \node[anchor=north west, red, scale=0.8, align=left] at (1.15,-1.05)
    {$C = 1-j$\\ $= \sqrt{2}e^{-j\pi/4}$};
\end{tikzpicture}
\end{document}
